\begin{lemma}
With \(s,m,p,k\) as in
\(\noderef{TwoBiteStretch}\), \(\noderef{TwoBiteReducedVertexCount}\),
\(\noderef{TwoBiteEdgeProbability}\), and
\(\noderef{TwoBiteIndependenceScale}\), and with
\(\beta=\frac12\) and \(\kappa=1+\varepsilon\), the parameters are positive for
all sufficiently large \(n\) and agree with the formulas used in the paper.
\end{lemma}
\begin{proof}
Fix \(\varepsilon,\beta,\kappa\) with \(0<\varepsilon\),
\(\beta=\frac12\), and \(\kappa=1+\varepsilon\).  Choose a threshold \(n_0\)
large enough that every \(n\ge n_0\) satisfies \(n>1\).  Then \(n>0\) and
\(\log n>0\).

For such an \(n\), unfold \(\noderef{TwoBiteStretch}\).  It gives
\[
s=\log^2 n.
\]
Since \(\log n>0\), its square is positive, so \(s>0\).  This is exactly the
first asserted positivity statement \(0<\noderef{TwoBiteStretch}(n)\).

Next unfold \(\noderef{TwoBiteEdgeProbability}\).  With \(\beta=\frac12\),
\[
p=\frac12\sqrt{\frac{\log n}{n}}.
\]
The factor \(\frac12\) is positive.  The quotient \(\frac{\log n}{n}\) is
positive because both numerator and denominator are positive.  Therefore its
square root is positive, and the product is positive.  Hence
\[
0<\noderef{TwoBiteEdgeProbability}\!\left(\frac12,n\right).
\]

Unfolding \(\noderef{TwoBiteReducedVertexCount}\) gives the reduced base-graph
size
\[
m=\frac{n}{\noderef{TwoBiteStretch}(n)}
  =\frac{n}{\log^2 n},
\]
using the already unfolded expression for \(s\).  This is the \(m=n/s\)
relation used in the paper.  Since \(n>0\) and \(s>0\), division by the
positive number \(s\) preserves positivity, so
\[
m=\frac ns>0.
\]

Finally unfold \(\noderef{TwoBiteIndependenceScale}\).  For every \(n\),
\[
\noderef{TwoBiteIndependenceScale}(\kappa,n)
  =\kappa\sqrt{n\log n}.
\]
Substituting \(\kappa=1+\varepsilon\) gives
\[
k=(1+\varepsilon)\sqrt{n\log n}.
\]
Moreover \(\kappa=1+\varepsilon>0\), because \(0<\varepsilon\).  Also
\(n\log n>0\), since both \(n\) and \(\log n\) are positive.  Hence
\(\sqrt{n\log n}>0\), and multiplying the two positive factors gives
\[
k=\kappa\sqrt{n\log n}>0.
\]

Thus, for all \(n\ge n_0\), all four parameters are positive and have the
displayed formulas:
\[
s=\log^2 n,\qquad
m=\frac{n}{s}=\frac{n}{\log^2 n},\qquad
p=\frac12\sqrt{\frac{\log n}{n}},\qquad
k=(1+\varepsilon)\sqrt{n\log n}.
\]
\end{proof}
